\documentclass[11pt]{article}
\usepackage[T1]{fontenc}
\usepackage[utf8]{inputenc}
\usepackage{lmodern}
\usepackage[margin=1in]{geometry}
\usepackage{amsmath,amssymb,amsthm,mathtools}
\usepackage{microtype,booktabs}
\usepackage[hidelinks]{hyperref}
\newtheorem{theorem}{Theorem}[section]
\newtheorem{proposition}[theorem]{Proposition}
\newtheorem{corollary}[theorem]{Corollary}
\newcommand{\A}{\mathbf A_I}
\newcommand{\Pow}{\mathcal P_+}
\newcommand{\ow}[1]{\mathbin{\star_{#1}}}
\newcommand{\at}[1]{\langle #1\rangle}
\newcommand{\All}{I}
\newcommand{\pv}{\mathsf p}
\newcommand{\nv}{\mathsf n}
\title{Small Matrices and Optimal Tests for Owned Choice}
\author{KFOL research working note}
\date{28 October 2026}
\begin{document}
\maketitle
\begin{abstract}
The intended signature includes internal role-switching; this note first
analyzes its reduct without that operation. The algebra of possible pure
subgame-perfect winners has $2^K-1$ values for $K\geq3$ roles. We compute
the reduct's exact contextual quotient for each
player's possible-win observation: it has $K+1$ values. These quotients
jointly give a faithful representation. They imply an all-arity test
theorem: a fixed input domain determines every term function exactly when
it contains every singleton and ambiguous inputs with empty intersection.
Its minimum size is six for three roles, and $K+2$ for $K\geq4$.
We give canonical multivariable forms, classify every congruence by a
closure operator, and obtain the subdirectly irreducible quotients and a
minimum subdirect decomposition. Adding internal role permutations
reduces the congruences to a chain of uniform-matroid truncations.
The observational congruence and the
general test-domain sufficiency theorem are kernel-checked in Lean.
Unary synthesis and the test-domain necessity theorem are now also
kernel-checked in \texttt{Synthesis.lean}.
The remaining proofs and their finite checks are identified explicitly.
\end{abstract}

\section{One small signature}

Fix a finite role set $I$, with $K=|I|\geq3$. The algebra studied here is
\[
 \A=(\Pow(I); (\star_r)_{r\in I},(\{i\})_{i\in I}),
 \qquad
 S\ow r T=\begin{cases}
 \{r\},&S=\{r\}\text{ or }T=\{r\},\\
 S\cup T,&\text{otherwise}.
 \end{cases}
\]
Its terms describe finite public games with role-owned binary choices and
one winner at each leaf. A player values its own win at one and all other
winners at zero. The denotation is the set of possible winners over full
pure subgame-perfect equilibrium (SPE) profiles. The earlier note,
\texttt{report.tex}, gives the full-strategy interpretation and its Lean
proof, the neutral-union construction, and the unary classification.

Ordinary union is a term operation: for distinct $r,s,t$,
\[
 (x\ow r y)\ow t(x\ow s y)=x\cup y.
\]
Every nonempty constant is consequently available. Adding a union symbol
is a finite definitional expansion of the signature. Each $\star_r$ is
associative, commutative, and idempotent, so a quantifier over a finite
nonempty domain is its fold in this fragment. This does not settle
infinitary or dependent quantification.

Role permutations are \emph{external symmetries} in this note. There is no
internal permutation constructor. The observation and test theorems below
refer to exactly this signature. An internal operation that transports an
input's role labels changes which one-player observations are available.

\section{Exactly what one player can observe}

Fix $k\in I$. An observer asks whether $k$ is a possible winner after a
term context has used the input. Put
\[
 S\equiv_k T\quad\Longleftrightarrow\quad
 \forall f\in\operatorname{Pol}_1(\A),\quad
       k\in f(S)\ \Longleftrightarrow\ k\in f(T).
\]
Here the polynomial may repeat the input arbitrarily and use all ground
constants. This is the standard contextual, or Leibniz, congruence of the
logical matrix whose designated values are the sets containing $k$
\cite{BlokPigozzi1989,Font1993}. Computing it is the specific problem.

\begin{theorem}[Exact one-player quotient]\label{thm:view}
For nonempty $S,T\subseteq I$,
\begin{equation}\label{eq:view}
 S\equiv_k T\quad\Longleftrightarrow\quad
 (k\in S\Longleftrightarrow k\in T)
 \ \land\ \bigwedge_{i\ne k}(S=\{i\}\Longleftrightarrow T=\{i\}).
\end{equation}
There are exactly $K+1$ classes:
\begin{enumerate}
\item $\pv$: every set containing $k$, including $\{k\}$;
\item $a_i$ for each $i\ne k$: the single value $\{i\}$;
\item $\nv$: every ambiguous set omitting $k$.
\end{enumerate}
This is the coarsest congruence preserving the possible-win observation
for $k$.
\end{theorem}
\begin{proof}
For nonempty $S,T$, direct calculation gives
\[
 S\ow r T=\{i\}\ \Longleftrightarrow\
 \begin{cases}
 S=\{i\}\text{ or }T=\{i\},&r=i,\\
 S=T=\{i\},&r\ne i.
 \end{cases}
\]
Also, membership of the owner $r$ is simply the disjunction of its
membership in the two inputs. For $k\ne r$, membership of $k$ is
\[
 (k\in S\text{ or }k\in T)\ \land\ S\ne\{r\}\ \land\ T\ne\{r\}.
\]
These equations show that the right-hand side of~\eqref{eq:view} is a
congruence, hence is preserved by every context.

Conversely, the identity context tests membership of $k$. For $i\ne k$,
the context $x\ow i\{k\}$ excludes $k$ exactly when $x=\{i\}$. These
contexts distinguish all the listed cases. The class $\nv$ is inhabited
because there are two distinct roles other than $k$.
\end{proof}

\texttt{sameView\_own} and \texttt{contextualAt\_iff} prove the
congruence and its contextual characterization generically in Lean.
Ground constants suffice to represent arbitrary polynomial parameters
because union realizes every nonempty subset.

Write $C_k=\A/{\equiv_k}$. Its operations have a short explicit form.
For owner $r\ne k$, an input $a_r$ absorbs. Without that absorber, an
input $\pv$ gives $\pv$; two equal atoms $a_i$ give $a_i$; all other
cases give $\nv$. For owner $k$, omit the absorbing-$a_r$ case.
Equivalently, each operation is the join of a simple order: distinct
non-owner atoms join to $\nv$, then comes $\pv$, and an owner atom
$a_r$, when present, sits above everything.

\begin{corollary}[Faithful family of small perspectives]\label{cor:product}
The canonical map
\[
 \A\longrightarrow\prod_{k\in I} C_k,\qquad
 S\longmapsto([S]_{\equiv_k})_{k\in I}
\]
is an injective homomorphism and is onto each factor. Thus the
$2^K-1$-valued algebra has a subdirect representation by $K$ factors,
each of size $K+1$.
\end{corollary}
\begin{proof}
Each coordinate is a quotient homomorphism. If two inputs differ, some
$k$ belongs to exactly one of them, and their $k$ coordinates differ.
\end{proof}

This retains possible SPE winners. It is not the Boolean guarantee
representation: $\{k\}$ and every ambiguous set containing $k$ are
\emph{identified} in $C_k$, whereas individual forceability distinguishes
them. The family of all perspectives jointly recovers every membership
bit. Role renaming transports $C_k$ to $C_{\pi k}$ without privileging
one observer.

It also supplies ordinary finite logical matrices: designate $\pv$ in
$C_k$ and interpret consequence by preservation of that designation.
Generating a truth-table proof system for a finite matrix is standard;
the content here is the explicit smallest contextual matrix.

\section{Optimal fixed input tests, at every arity}

Call $D\subseteq\Pow(I)$ a complete input test domain if, for every arity
$n$ and every pair of terms $f,g$,
\[
 f=g\text{ on }\Pow(I)^n\quad\Longleftrightarrow\quad f=g\text{ on }D^n.
\]
The domain is chosen once and must work for all terms and all arities.

\begin{theorem}[Exact characterization]\label{thm:tests}
A domain $D$ is complete exactly when:
\begin{enumerate}
\item every singleton $\{i\}$ belongs to $D$;
\item for every $k$, some ambiguous $F\in D$ omits $k$.
\end{enumerate}
The second condition says that the intersection of the ambiguous members
of $D$ is empty, with the intersection of an empty family taken as $I$.
\end{theorem}
\begin{proof}
\emph{Sufficiency.} Suppose $f,g$ differ at a valuation $v$, and choose
$k$ belonging to exactly one output. Choose an ambiguous $F\in D$
omitting $k$. Replace each input set $S$ as follows:
\[
 S\longmapsto
 \begin{cases}
 \{k\},&k\in S,\\
 \{i\},&S=\{i\},\ i\ne k,\\
 F,&\text{otherwise}.
 \end{cases}
\]
The replacement belongs to $D$ and has the same $k$-view. All operations
preserve that view, so the outputs still differ at $k$. The replacement
may use the \emph{same} $F$ for every variable. This proof works at any
arity, not just for unary functions.

\emph{Necessity of singleton tests.} The unary classification from the
earlier note realizes all functions
\[
 R(C,E)(S)=\begin{cases}\{i\}\cup E_i,&S=\{i\},\\
 S\cup C,&|S|\geq2,\end{cases}
 \quad E_i\subseteq C\setminus\{i\}.
\]
Fix $i$ and $k\ne i$. Take $C=I$, all rows full except
$E_i=I\setminus\{i,k\}$. This function differs from the constant $I$
only at $\{i\}$. Hence that singleton must be tested.

\emph{Necessity of ambiguity coverage.} For each $k$, the same
classification realizes
\[
 F_k(S)=\begin{cases}S,&|S|=1,\\S\cup\{k\},&|S|\geq2.\end{cases}
\]
This differs from the identity exactly on ambiguous sets omitting $k$.
At least one such input must be tested.
\end{proof}

The sufficiency statement is \texttt{complete\_test\_domain} in Lean.
The necessity statement is now \texttt{test\_domain\_necessary} in
\texttt{Synthesis.lean}, using explicit unary synthesis and term
separators. The cardinality corollary below remains a written argument.

\begin{corollary}[Sharp size]\label{cor:optimal}
The minimum size of a complete fixed test domain is
\[
 \begin{cases}6,&K=3,\\K+2,&K\geq4.\end{cases}
\]
For three roles, take every singleton and every pair. For four or more,
take every singleton and any two disjoint pairs.
\end{corollary}
\begin{proof}
Every singleton is necessary. One ambiguous set has nonempty
intersection with itself, so at least two are needed. Two disjoint pairs
suffice for $K\geq4$. At $K=3$, every ambiguous set has size at least
two, so any two intersect. All three pairs are necessary and sufficient.
\end{proof}

For a fixed observed player, an even smaller tailored domain consists of
all $K$ singletons and one fixed pair omitting that player. Thus an
equation can be checked in $K$ quotient algebras on $K+1$ states, or on a
single fixed input alphabet of the optimal size above. The latter input
alphabet is not a subalgebra: intermediate outputs may be larger sets.
No claim is made that it is a smaller carrier with the same operations.

\section{Canonical forms for several variables}

The small quotients already give compositional function tables at all
arities. There is a second, more syntactic normalization.

Fix a status assignment
$\sigma\in(I\cup\{*\})^n$, where status $i$ means the input is exactly
$\{i\}$, and $*$ means it is ambiguous. On this status cell, represent an
expression by $(C,J)$, denoting
\[
 C\ \cup\ \bigcup_{j\in J}S_j,
 \qquad J\subseteq\{j:\sigma_j=*\}.
\]
For a constant use $(\{i\},\varnothing)$; for a singleton-status variable
use that singleton; for an ambiguous variable use
$(\varnothing,\{j\})$. An owned node collapses to
$(\{r\},\varnothing)$ if either child has that form. Otherwise union the
two $C$ fields and the two $J$ fields. Finally, when $C=I$, discard $J$.

\begin{theorem}[Canonical status forms]\label{thm:normal}
This normalization is sound on each status cell. Its final pair $(C,J)$
is uniquely determined by the term function on that cell. Consequently
the table of these pairs over all $(K+1)^n$ statuses is a canonical
representation of any $n$-variable term function.
\end{theorem}
\begin{proof}
If $J$ is nonempty, the value contains an ambiguous input, so cannot be
an owner singleton. If $J$ is empty, the value is exactly $C$. The node
rule therefore implements precisely the two cases of $\star_r$.

For uniqueness, if two $C$ fields differ, choose $k$ in their symmetric
difference and give every ambiguous variable a pair omitting $k$.
Output membership of $k$ then detects the different fields. If $C$
agrees and differs from $I$, but $J$ differs, choose $k\notin C$, give one
variable in the symmetric difference a pair containing $k$, and give all
other ambiguous variables pairs omitting $k$. This separates the
outputs. The only redundant case is $C=I$, already canonicalized by
discarding $J$. All status cells are inhabited for $K\geq3$.
\end{proof}

These rules give a sound and complete normalization decision procedure for term
equality: compute both tables and compare them. They also explain how
ambiguity survives. Once the singleton cases at a cell are fixed, the
remaining computation is just union with a record of which input
variables contribute. The table is not a characterization of which
arbitrary tables are definable by terms, and hence does not yet classify
the full clone. Nor is this a finite unconditional equational basis.

\texttt{structure.py} implements both representations and a complete
counterexample search over the optimal input domain. The status-form
soundness and uniqueness arguments above are not yet mechanized.

\section{All congruences and the small irreducible quotients}

Union is term definable, so every algebra congruence is also a union
congruence. The standard correspondence with closure operators is useful
here; the class-maximum construction is classical, compare the lemma to
Proposition 6 in Crapo\cite{Crapo1967}. We give the finite argument
explicitly because the extra owner operations impose a strong condition.

\begin{theorem}[Congruence classification]\label{thm:congruences}
The congruences of $\A$ are exactly the kernels of closure operators
$c:\Pow(I)\to\Pow(I)$ satisfying
\begin{equation}\label{eq:atomclosure}
 c(\{i\})\in\{\{i\},I\}\qquad(i\in I).
\end{equation}
Equivalently, they correspond to families $\mathcal F\subseteq\Pow(I)$
such that $I\in\mathcal F$, $\mathcal F$ is closed under nonempty
intersections when the result is nonempty, and
\[
 S\in\mathcal F,\ S\ne I,\ i\in S\quad\Longrightarrow\quad
 \{i\}\in\mathcal F.
\]
The closure is the intersection of all members containing the input.
\end{theorem}
\begin{proof}
For a congruence $\theta$, define
$c(S)=\bigcup[S]_\theta$. Finiteness and union compatibility make this
maximum a member of its class. Extensivity and idempotence follow.
If $S\subseteq T$, union compatibility gives
$T\mathrel\theta c(S)\cup T$, so $c(S)\subseteq c(T)$. Thus $c$ is a
closure operator and $\theta=\ker c$.

If $\{i\}$ is identified with any different value $S$, apply
$x\mapsto x\ow i I$. It sends $\{i\}$ to $\{i\}$ and every other input
to $I$. Hence $\{i\}\mathrel\theta I$ and $c(\{i\})=I$.
This crucial step is \texttt{atom\_collapse\_to\_top} in Lean.

Conversely, any closure operator has a union-compatible kernel, since
$c(S\cup T)=c(c(S)\cup c(T))$. Assume~\eqref{eq:atomclosure}. Compare
equivalent input pairs for $\star_i$. If neither side absorbs, their
outputs are unions and have equal closures. If both absorb, both outputs
are $\{i\}$. If just one absorbs, an input equivalent to $\{i\}$ differs
from it. Their common closure cannot be $\{i\}$, since extensivity would
make that nonempty input equal to $\{i\}$. Their closure is therefore
$I$, and the nonabsorbing side's union also closes to $I$. This proves compatibility
with every owned operation.

The usual fixed-set description of closure operators gives the final
formulation. A proper closed set containing $i$ rules out
$c(\{i\})=I$ and therefore forces $\{i\}$ to be closed. Conversely, that
condition makes every nonclosed singleton close to $I$.
\end{proof}

In particular, there are exactly $18$ congruences at $K=3$ and $626$ at
$K=4$. The script independently enumerates all partitions at $K=3$.
At $K=4$ it enumerates the classified closure families and checks every
constructed kernel against every operation; completeness uses the proof
above, rather than a partition search.

\begin{theorem}[Subdirectly irreducible homomorphic images]\label{thm:si}
Every nontrivial subdirectly irreducible homomorphic image of $\A$ is
obtained from a nonempty proper set $S\subset I$ with closed family
\[
 \mathcal F_S=\{I,S\}\cup\{\{i\}:i\in S\}.
\]
For $|S|=1$ the quotient has two elements. For $|S|\geq2$ it has
$|S|+2$ elements: one singleton state per $i\in S$, the state $S$, and
the state $I$. In particular, every such image has at most $K+1$ elements.
The case $S=I\setminus\{k\}$ is exactly $C_k$.
\end{theorem}
\begin{proof}
The congruence/fixed-family correspondence reverses inclusion.
Meet-irreducible congruences therefore correspond to join-irreducible
admissible fixed families. Every admissible family is the join of the
families $\mathcal F_S$ for its proper members $S$. Each
$\mathcal F_S$ is join-irreducible: its member $S$ cannot be generated by
intersections of its strictly smaller members, which are singletons.
Thus these are exactly the meet-irreducible kernels and hence exactly the
nontrivial subdirectly irreducible quotients.

The displayed fixed sets are representatives of all quotient classes.
For $|S|\geq2$, its least nontrivial congruence identifies the states $S$
and $I$. To see minimality directly, collapsing an atom with $S$ forces
it to collapse with $I$ by its owned operation. Collapsing an atom with
$I$ then collapses $S$ with $I$ by combining it with a different atom
under an owner outside $S$. Collapsing two atoms first collapses one
with $I$ by its owned operation. Every nontrivial collapse therefore
contains the stated pair.
\end{proof}

\begin{corollary}[Minimum subdirect decomposition]
Every subdirect representation of $\A$ by nontrivial subdirectly
irreducible factors must include, up to the quotient maps' kernel,
each $C_k$. The $K$-factor decomposition of
Corollary~\ref{cor:product} is consequently minimum and irredundant.
\end{corollary}
\begin{proof}
The distinct values $I$ and $I\setminus\{k\}$ must be separated by some
factor. Among the classified quotients, only
$\mathcal F_{I\setminus\{k\}}$ separates them: every other proper
fixed set fails to contain $I\setminus\{k\}$, so that value closes to
$I$. Thus the corresponding factor $C_k$ is necessary for every $k$.
\end{proof}

These statements concern homomorphic images of $\A$ itself. They do not
bound subdirectly irreducible members of the whole variety generated by
$\A$, and do not establish finite basability.

\section{What internal role switching changes}

Now expand the signature by unary operations $S\mapsto\pi S$ for every
role permutation $\pi$. Write $\A^{\mathrm{sym}}$ for this algebra. This
is a separate extension: the preceding small test-domain and one-player
quotient theorems used the original signature.

\begin{theorem}[The symmetric congruence chain]\label{thm:sym}
There are exactly $K$ congruences of $\A^{\mathrm{sym}}$. They are the
kernels of
\[
 c_m(S)=\begin{cases}S,&|S|\leq m,\\I,&|S|>m,\end{cases}
 \qquad m=0,1,\ldots,K-1.
\]
Here $m=0$ gives the universal relation and $m=K-1$ gives equality.
For $1\leq m\leq K-1$, the quotient has
$1+\sum_{j=1}^{m}\binom Kj$ values. These are the restrictions to
nonempty sets of the uniform-matroid closures of rank $m+1$.
\end{theorem}
\begin{proof}
A congruence of the expansion is exactly a congruence of $\A$ invariant
under every permutation. Its class-maximum closure therefore commutes
with permutations. If any singleton closes to $I$, symmetry makes every
singleton close to $I$, and monotonicity makes every nonempty set close
to $I$.

Otherwise all singletons are closed. Let $m$ be the largest size of a
proper closed set. Permutation invariance makes all $m$-element sets
closed. Their nonempty intersections give every nonempty set of size at
most $m$: for a smaller set $T$ and $x\notin T$, an $m$-element superset
of $T$ omitting $x$ exists because $m<K$. Thus the proper closed sets
are exactly those of size at most $m$. The congruence classification
gives precisely the displayed kernels. Their closures are nested, so
their kernels form a chain. The cardinalities and the uniform-matroid
identification follow from their fixed sets; see the uniform matroids
in Oxley\cite{Oxley2011}.
\end{proof}

\begin{corollary}[Symmetry makes the algebra irreducible]
$\A^{\mathrm{sym}}$ is subdirectly irreducible. Its least nontrivial
congruence identifies all sets of size at least $K-1$, leaving every
smaller nonempty set fixed.
\end{corollary}
\begin{proof}
In the chain above, $\ker c_{K-2}$ is the least congruence strictly above
equality. It exists for $K\geq3$ and has the stated classes.
\end{proof}

The quotient with $m=1$ records every unique winner separately and merges
all ambiguity. The intervening quotients retain increasingly large sets
of alternative equilibrium winners. Thus internal role symmetry permits
an exact hierarchy of ambiguity thresholds.

The one-player contextual quotient also changes: it is now equality.
If $S,T$ differ at role $i$, choose a permutation taking $i$ to $k$.
The context $\pi(x)$ then distinguishes their possible-win observation
for $k$. This does not invalidate the earlier product representation:
the tuple of all $C_k$ still represents the underlying algebra, and
permutations act by transporting \emph{between} its coordinates. They
cannot be interpreted as separate unary operations inside each fixed
factor. The distinction between external symmetry and internal role
switching therefore has a measurable algebraic effect.

\section{Repetition and substitution remain separate questions}

For ordinary unary terms $f$, define uniform-profile observation by
\[
 U(f)(S)=\bigcup_{i\in S}f(\{i\}).
\]
Fixing the same full component SPE profile in every occurrence gives
exactly this function: the surrounding tree sees its winner as a leaf.
The uniform-copy result for neutral union is a special positive case.

The earlier note's term
$D(x)=x\ow c((x\ow a\{b\})\ow b\{a\})$ fixes all singleton inputs but
adds $\{a,b\}$ to ambiguous inputs when $I=\{a,b,c\}$. Thus $U(D)$ is
the identity. Nevertheless, for $H(x)=x\cup\{c\}$,
\[
 U(D\circ H)(\{a\})=I,\qquad
 (U(D)\circ U(H))(\{a\})=\{a,c\}.
\]
The uniform-profile equivalence is not preserved by substituting an
expression for its variable. Binding the entire component's selected
behavior differs from sharing only the input behavior inside several
separately solved surrounding contexts. A typed treatment of sharing
must identify that boundary explicitly. These observations are about
ordinary SPE profiles selected for copy equality; they do not restrict
deviations or define an imperfect-information equilibrium.

\begin{proposition}[Exactly which unary contexts tolerate uniform selection]
For $R(C,E)$ from the unary classification, $U(R)=R$ exactly when, for
every $k\in C$, at most one $i\ne k$ has $k\notin E_i$. There are exactly
\[
 (K+1)^K+2^K-2
\]
ordinary unary term functions preserved by $U$, whereas the number of
distinct functions in the image of $U$ is
\[
 2^{K(K-1)}+2^K-2.
\]
For $K=3$ the two counts are both $70$; for $K=4$ they are $639$ and
$4110$, respectively.
\end{proposition}
\begin{proof}
On $S$, uniform lifting returns
$S\cup\bigcup_{i\in S}E_i$. For each extra winner $k\in C$, every
ambiguous input omitting $k$ must therefore meet the rows permitting
$k$. This holds exactly when no pair of other rows both omit $k$.
For each $k$, either it is absent from $C$, or it is present and all
other rows permit it except at most one. There are $1+K$ options per
$k$. Add the nonempty constant functions and subtract the overlap $I$.

For the image count, any row family $E_i\subseteq I\setminus\{i\}$ is
realizable by setting $C=I$ in the earlier synthesis. Singleton rows
uniquely determine their union-preserving lifts. There are
$2^{K(K-1)}$ such extensive lifts. Adding constants again gives the
displayed count.
\end{proof}

For $K\geq4$, uniform lifting can even leave the ordinary unary term
functions. Choose distinct $a,b,c,d$ and a regular function with
$C=\{d\}$, $E_a=\{d\}$, and every other row empty. Its lift adds $d$
exactly when the input contains $a$. It adds $d$ at $\{a,b\}$ but not
at $\{b,c\}$, although both are ambiguous and omit $d$. No regular
ordinary term function can do that: it would need $d$ both present and
absent in its fixed extra set. This further separates the two accounts
of repetition without introducing beliefs or probabilities.

\section{What is established and what remains}

\begin{center}\small
\begin{tabular}{@{}p{0.47\linewidth}p{0.45\linewidth}@{}}
\toprule
Statement & Evidence\\\midrule
Exact contextual relation and congruence & Generic Lean proof\\
Arbitrary-arity test-domain sufficiency & Generic Lean proof\\
Optimal domain characterization and size & Generic Lean proofs of both conditions; cardinality argument on paper; finite separation checks\\
Canonical status forms & Mathematical proof and finite binary/ternary cross-checks\\
All congruences and SI quotient classification & Mathematical proof; atom-collapse lemma in Lean; finite congruence checks\\
Minimum subdirect decomposition & Mathematical consequence of the classifications\\
Internal-permutation congruence chain & Mathematical proof; exhaustive invariance tests among the classified congruences at $K=3,4$\\
Priority of these specific classifications & Not established\\\bottomrule
\end{tabular}
\end{center}

The generic quotient method is established algebraic logic
\cite{Font1993}; the closure-operator construction is established
semilattice theory~\cite{Crapo1967}. We reuse both. The concrete
calculations for the SPE-winner algebra, optimal all-arity test domains,
and its irreducible quotient structure are the candidate contributions.
Targeted searches did not locate these statements, which is not a
priority proof. The known owner-indexed and guarantee-based precedents
remain essential comparisons~\cite{Abramsky2007,OldeLoohuisVenema2010}.

The next algebraic question is whether the canonical normalization can
be replaced by a useful finite unconditional equational basis, or
whether the family admits a concise uniform presentation as $K$ varies.
The current normalization already handles several variables and gives
complete finite counterexample search. It does not require adding hidden
information, probabilities, or beliefs to the syntax.

\bibliographystyle{plain}
\bibliography{references}
\end{document}
